% GENERATED FILE. Edit canonical lessons, then run npm run generate:books.
% canonical-source-hash: fnv1a64:edae598106cbe2c0
% canonical-lessons: GE-C418-flugbegleiter, GE-C418-fensterplatz, GE-C418-gangplatz, GE-C418-abfliegen, GE-C418-landung

\chapter{An Bord}
\label{ch:ge-an-bord}
\hlchaptermodality{418}

\begin{chapteropening}
\textbf{By the end of this chapter:} \emph{I can talk about a flight: the flight attendant, the window seat, the aisle seat, taking off and the landing.}

Complete the last lesson of chapter 418: I can talk about a flight: the flight attendant, the window seat, the aisle seat, taking off and the landing.
\end{chapteropening}

\section[der Flugbegleiter]{der Flugbegleiter — the flight attendant}

\label{lesson:GE-C418-flugbegleiter}



\begin{quote}
Before the new one: say the German for to declare, then the German for the passport.
\end{quote}

\subsection*{You'll want to know: der Flugbegleiter}
\textbf{der Flugbegleiter} — \textquotedblleft{}the flight attendant\textquotedblright{}.

\begin{quote}
Der Flugbegleiter bringt uns Wasser und Kaffee.
\end{quote}

\begin{quote}
\emph{The flight attendant brings us water and coffee.}
\end{quote}

Same form in the plural. A woman: \emph{die Flugbegleiterin}.

\begin{cousinweb}
\emph{Flugbegleiter} is built from \emph{Flug} 'flight', inherited from Proto-Germanic \emph{flugiz} (EtymDB), and \emph{Begleiter} 'companion', from \emph{begleiten} 'to accompany', which is built on \emph{leiten} 'to lead'.
\end{cousinweb}

\begin{grammarlens}[title={What you've built}]
One more word for telling, explaining and describing at B1.
\end{grammarlens}

\subsection*{Your turn}
\begin{itemize}
\raggedright
  \item \emph{Say it:} \emph{der Flugbegleiter}
  \item \emph{Say it:} \emph{der Flugbegleiter}, once more
  \item \emph{Say it:} a sentence of your own with \emph{der Flugbegleiter}
  \item \emph{Recall:} say the German for to declare, then the German for the passport, then say \emph{der Flugbegleiter} again
\end{itemize}

\begin{tcolorbox}[breakable,colback=teal!4,colframe=teal!35!black,title={Before you move on}]
What does der Flugbegleiter mean? (\textquotedblleft{}The flight attendant\textquotedblright{}.) Say it once more.
\end{tcolorbox}
\section[der Fensterplatz]{der Fensterplatz — the window seat}

\label{lesson:GE-C418-fensterplatz}



\begin{quote}
Before the new one: say the German for the passport, then the German for the flight attendant.
\end{quote}

\subsection*{You'll want to know: der Fensterplatz}
\textbf{der Fensterplatz} — \textquotedblleft{}the window seat\textquotedblright{}.

\begin{quote}
Mein Sohn möchte immer einen Fensterplatz.
\end{quote}

\begin{quote}
\emph{My son always wants a window seat.}
\end{quote}

Plural: \emph{die Fensterplätze}.

\begin{cousinweb}
\emph{Fensterplatz} is built from \emph{Fenster} 'window', an early borrowing from Latin \emph{fenestra} (Lewis \& Short), and \emph{Platz} 'place', which came through Old French \emph{place} from Latin \emph{platea} 'broad street' (EtymDB).
\end{cousinweb}

\begin{grammarlens}[title={What you've built}]
One more word for telling, explaining and describing at B1.
\end{grammarlens}

\subsection*{Your turn}
\begin{itemize}
\raggedright
  \item \emph{Say it:} \emph{der Fensterplatz}
  \item \emph{Say it:} \emph{der Fensterplatz}, once more
  \item \emph{Say it:} a sentence of your own with \emph{der Fensterplatz}
  \item \emph{Recall:} say the German for the passport, then the German for the flight attendant, then say \emph{der Fensterplatz} again
\end{itemize}

\begin{tcolorbox}[breakable,colback=teal!4,colframe=teal!35!black,title={Before you move on}]
What does der Fensterplatz mean? (\textquotedblleft{}The window seat\textquotedblright{}.) Say it once more.
\end{tcolorbox}
\section[der Gangplatz]{der Gangplatz — the aisle seat}

\label{lesson:GE-C418-gangplatz}



\begin{quote}
Before the new one: say the German for the flight attendant, then the German for the window seat.
\end{quote}

\subsection*{You'll want to know: der Gangplatz}
\textbf{der Gangplatz} — \textquotedblleft{}the aisle seat\textquotedblright{}.

\begin{quote}
Auf langen Flügen nimmt er lieber einen Gangplatz.
\end{quote}

\begin{quote}
\emph{On long flights he prefers an aisle seat.}
\end{quote}

Plural: \emph{die Gangplätze}. \emph{Der Gang} — the aisle, the corridor.

\begin{cousinweb}
\emph{Gangplatz} is built from \emph{Gang} 'passage, aisle', from \emph{gehen} 'to go', inherited from Old High German \emph{gang} (EtymDB), and \emph{Platz} 'place', from Latin \emph{platea}.
\end{cousinweb}

\begin{grammarlens}[title={What you've built}]
One more word for telling, explaining and describing at B1.
\end{grammarlens}

\subsection*{Your turn}
\begin{itemize}
\raggedright
  \item \emph{Say it:} \emph{der Gangplatz}
  \item \emph{Say it:} \emph{der Gangplatz}, once more
  \item \emph{Say it:} a sentence of your own with \emph{der Gangplatz}
  \item \emph{Recall:} say the German for the flight attendant, then the German for the window seat, then say \emph{der Gangplatz} again
\end{itemize}

\begin{tcolorbox}[breakable,colback=teal!4,colframe=teal!35!black,title={Before you move on}]
What does der Gangplatz mean? (\textquotedblleft{}The aisle seat\textquotedblright{}.) Say it once more.
\end{tcolorbox}
\section[abfliegen]{abfliegen — to take off, to fly out}

\label{lesson:GE-C418-abfliegen}



\begin{quote}
Before the new one: say the German for the window seat, then the German for the aisle seat.
\end{quote}

\subsection*{You'll want to know: abfliegen}
\textbf{abfliegen} — \textquotedblleft{}to take off, to fly out\textquotedblright{}.

\begin{quote}
Das Flugzeug fliegt pünktlich um sieben ab.
\end{quote}

\begin{quote}
\emph{The plane takes off on time at seven.}
\end{quote}

Separable; Perfekt with \emph{sein}: \emph{wir sind abgeflogen}. \emph{Der Abflug} — the departure.

\begin{cousinweb}
\emph{abfliegen} is built from \emph{ab} 'off, away' and \emph{fliegen} 'to fly', inherited from Proto-Germanic \emph{fleuganą} (EtymDB). English \emph{fly} is a cousin.
\end{cousinweb}

\begin{grammarlens}[title={What you've built}]
One more word for telling, explaining and describing at B1.
\end{grammarlens}

\subsection*{Your turn}
\begin{itemize}
\raggedright
  \item \emph{Say it:} \emph{abfliegen}
  \item \emph{Say it:} \emph{abfliegen}, once more
  \item \emph{Say it:} a sentence of your own with \emph{abfliegen}
  \item \emph{Recall:} say the German for the window seat, then the German for the aisle seat, then say \emph{abfliegen} again
\end{itemize}

\begin{tcolorbox}[breakable,colback=teal!4,colframe=teal!35!black,title={Before you move on}]
What does abfliegen mean? (\textquotedblleft{}To take off, to fly out\textquotedblright{}.) Say it once more.
\end{tcolorbox}
\section[die Landung]{die Landung — the landing}

\label{lesson:GE-C418-landung}



\begin{quote}
Before the new one: say the German for the aisle seat, then the German for to take off.
\end{quote}

\subsection*{You'll want to know: die Landung}
\textbf{die Landung} — \textquotedblleft{}the landing\textquotedblright{}.

\begin{quote}
Die Landung war ruhig und pünktlich.
\end{quote}

\begin{quote}
\emph{The landing was smooth and on time.}
\end{quote}

Plural: \emph{die Landungen}. \emph{Landen} — to land.

\begin{cousinweb}
\emph{Landung} is built on \emph{landen} 'to land', from \emph{Land} 'land', inherited from Proto-Germanic \emph{landą} (EtymDB).
\end{cousinweb}

\begin{grammarlens}[title={What you've built}]
One more word for telling, explaining and describing at B1.
\end{grammarlens}

\subsection*{Your turn}
\begin{itemize}
\raggedright
  \item \emph{Say it:} \emph{die Landung}
  \item \emph{Say it:} \emph{die Landung}, once more
  \item \emph{Say it:} a sentence of your own with \emph{die Landung}
  \item \emph{Recall:} say the German for the aisle seat, then the German for to take off, then say \emph{die Landung} again
\end{itemize}

\begin{tcolorbox}[breakable,colback=teal!4,colframe=teal!35!black,title={Before you move on}]
What does die Landung mean? (\textquotedblleft{}The landing\textquotedblright{}.) Say it once more.
\end{tcolorbox}
